\section{Circuit Complexity --- DAGTransform}
%%% generated by the blueprint pipeline from TCSlib.Complexity.CircuitComplexity.DAGTransform
\subsection*{Circuit Complexity --- DAGTransform: Introduction}

Declarations of the Lean module \texttt{TCSlib.Complexity.CircuitComplexity.DAGTransform}.

\subsection*{Circuit Complexity --- DAGTransform: Results}

\begin{definition}[Specification of a gate gadget]
\lean{BoolCircuit.GadgetSpec}
\label{BoolCircuit.GadgetSpec}
\leanok
\uses{BoolCircuit.DAGGate, BoolCircuit.DAGGate.eval, BoolCircuit.GatesAcyclic, BoolCircuit.runWith, BoolCircuit.vertexDepth, BoolCircuit.vertexValue}
Let $n,c,	ext{cost}$ be natural numbers, let $P$ be a property of gates, let $g$ be a
gate, let $gs$ be a list of gates, and let $r=(r_1,r_2)$ be a list of gates together
with a natural-number vertex. A gadget specification requires that $r_1$ extend $gs$ by
appending some list; that $r_1$ be acyclic on $n$ inputs; that $r_2<n+|r_1|$; and that,
on every Boolean input, the value at $r_2$ in $r_1$ equal the evaluation of $g$ on the
values obtained by running $gs$ on that input. It also requires $|r_1|\le|gs|+	ext{cost}$,
depth at $r_2$ at most $c$ plus the maximum depth in $gs$ of the arguments of $g$ (with
empty maximum taken as $0$), and that every gate appended to $gs$ satisfy $P$.
\end{definition}

\begin{definition}[Correctness of a gate gadget]
\lean{BoolCircuit.GadgetCorrect}
\label{BoolCircuit.GadgetCorrect}
\leanok
\uses{BoolCircuit.DAGGate, BoolCircuit.GadgetSpec, BoolCircuit.GateKind, BoolCircuit.GatesAcyclic}
For natural numbers $n,c$, a predicate $Q$ on gate kinds and input arities, a predicate
$P$ on gates, and an emitter that assigns an output vertex and a list of gates to each
gate and existing gate list, gadget correctness means the following: for every gate $g$
and gate list $G$ that is acyclic on $n$ inputs, if every argument of $g$ is a vertex
among the $n$ inputs and the vertices of $G$, and $Q$ holds for the kind of $g$ and its
number of arguments, then the emitter's result satisfies the gadget specification with
depth overhead $c$, size allowance $\operatorname{arity}(g)+2$, and new-gate predicate
$P$.
\end{definition}

\begin{definition}[One circuit rewriting step]
\lean{BoolCircuit.rewriteStep}
\label{BoolCircuit.rewriteStep}
\leanok
\uses{BoolCircuit.DAGGate}
Given an emitter, a state consisting of a list of rewritten gates and a list of vertex
indices, and a gate $g$, replace each argument index of $g$ by the corresponding entry
of the index list, using $0$ when the index is out of range. Apply the emitter to this
relabeled gate and the rewritten gate list, then keep the emitter's new gate list and
append its output index to the index list.
\end{definition}

\begin{definition}[Rewriting a list of gates]
\lean{BoolCircuit.rewriteGates}
\label{BoolCircuit.rewriteGates}
\leanok
\uses{BoolCircuit.DAGGate, BoolCircuit.rewriteStep}
Given a natural number $n$, an emitter, and a list of gates, process the gates from left
to right by the one-step rewriting procedure, starting with an empty rewritten gate list
and the index list $[0,1,\ldots,n-1]$. Return the resulting rewritten gate list and
index list.
\end{definition}

\begin{definition}[Invariants of a rewritten circuit]
\lean{BoolCircuit.RewriteInv}
\label{BoolCircuit.RewriteInv}
\leanok
\uses{BoolCircuit.DAGGate, BoolCircuit.DAGGate.depth, BoolCircuit.DAGGate.eval, BoolCircuit.GatesAcyclic, BoolCircuit.runWith, BoolCircuit.vertexDepth, BoolCircuit.vertexValue}
For natural numbers $n,c$, a gate predicate $P$, an original gate list $G$, and a state
$(H,m)$ consisting of rewritten gates and a vertex-index list, the rewrite invariant
requires all of the following: $m$ has length $n+|G|$; each entry of $m$ at an index
below $n+|G|$ is a valid vertex of $H$; $H$ is acyclic on $n$ inputs; and, for every
Boolean assignment to those inputs, the value at each such mapped vertex equals the
value at the corresponding vertex after evaluating $G$. It also requires that each
mapped vertex's depth in $H$ is at most $c$ times the corresponding original depth, that
$|H|$ is at most the sum over gates of their arities plus two, and that every gate in
$H$ satisfies $P$.
\end{definition}

\begin{theorem}[Bound on a list maximum]
\lean{BoolCircuit.foldr_max_le'}
\label{BoolCircuit.foldr_max_le'}
\leanok
\difficulty{4}
Let $l$ be a finite list of natural numbers, let $f:\bbr	o\bbr$ be a function on the
natural numbers, and let $B$ be a natural number. If $f(a)\le B$ for every $a$ in $l$,
then the maximum of the values $f(a)$ over $a$ in $l$, taken to be $0$ when $l$ is
empty, is at most $B$.
\end{theorem}

\begin{theorem}[Bound by the maximum of mapped list values]
\lean{BoolCircuit.le_foldr_max'}
\label{BoolCircuit.le_foldr_max'}
\leanok
\difficulty{4}
Let $l$ be a list of natural numbers, let $f:\mathbb{N}\to\mathbb{N}$, and let $a\in l$.
Then $f(a)$ is at most the maximum of the values of $f$ on the entries of $l$, with the
maximum taken together with $0$.
\end{theorem}

\begin{theorem}[Invariant for the empty gate list]
\lean{BoolCircuit.rewriteInv_nil}
\label{BoolCircuit.rewriteInv_nil}
\leanok
\uses{BoolCircuit.DAGGate, BoolCircuit.GatesAcyclic.nil, BoolCircuit.RewriteInv, BoolCircuit.vertexDepth, BoolCircuit.vertexDepth_input, BoolCircuit.vertexValue, BoolCircuit.vertexValue_input}
\difficulty{4}
For any natural numbers $n$ and $c$ and any gate property $P$, the empty gate list with
the identity map on the $n$ input vertices satisfies the rewriting invariant: the map
has length $n$, each input maps to a valid vertex, the rewritten list is acyclic and has
length zero, and the mapped values and depths satisfy the invariant bounds; every gate
in the rewritten list satisfies $P$.
\end{theorem}

\begin{theorem}[Preservation of the invariant in one rewrite step]
\lean{BoolCircuit.rewriteInv_step}
\label{BoolCircuit.rewriteInv_step}
\leanok
\uses{BoolCircuit.DAGGate, BoolCircuit.DAGGate.depth, BoolCircuit.DAGGate.eval, BoolCircuit.DAGGate.eval_remap, BoolCircuit.GadgetCorrect, BoolCircuit.GateKind, BoolCircuit.GatesAcyclic, BoolCircuit.RewriteInv, BoolCircuit.foldr_max_le', BoolCircuit.le_foldr_max', BoolCircuit.rewriteStep, BoolCircuit.runWith, BoolCircuit.runWith_append, BoolCircuit.runWith_getD_last, BoolCircuit.runWith_getD_of_lt, BoolCircuit.vertexDepth, BoolCircuit.vertexDepth_append, BoolCircuit.vertexValue_append}
\difficulty{6}
Suppose a gate-emission procedure satisfies the gadget specification for every eligible
gate and acyclic existing list. If a rewritten state satisfies the invariant for an
acyclic original list, adjoining a gate $g$ preserves the invariant provided the
enlarged original list remains acyclic and the gate's kind and argument count satisfy
$Q$. The resulting state is obtained by one rewrite step using the emission procedure.
\end{theorem}

\begin{theorem}[Invariant under gate rewriting]
\lean{BoolCircuit.rewriteGates_inv}
\label{BoolCircuit.rewriteGates_inv}
\leanok
\uses{BoolCircuit.DAGGate, BoolCircuit.GadgetCorrect, BoolCircuit.GateKind, BoolCircuit.GatesAcyclic, BoolCircuit.RewriteInv, BoolCircuit.rewriteGates, BoolCircuit.rewriteInv_nil, BoolCircuit.rewriteInv_step}
\difficulty{4}
Suppose a gate-emission procedure is correct in the following sense: for every acyclic
existing gate list, every gate whose arguments refer only to available vertices, and
every gate kind and argument count satisfying a predicate $Q$, the emitted gadget
extends the list and satisfies the stated gadget specification with depth factor $c$,
cost equal to the gate's argument count plus $2$, and new-gate property $P$. Then for
every acyclic original gate list whose gates satisfy $Q$, rewriting the list preserves
the invariant that the rewritten list is acyclic, has length at most the sum of each
original gate's argument count plus $2$, and contains only gates satisfying $P$; its
vertex map has length $n$ plus the number of original gates and maps each original
vertex to a valid rewritten vertex, preserving all vertex values and bounding each
mapped depth by $c$ times the original depth.
\end{theorem}

\begin{definition}[Rewriting a directed acyclic circuit]
\lean{BoolCircuit.DAGCircuit.rewrite}
\label{BoolCircuit.DAGCircuit.rewrite}
\leanok
\uses{BoolCircuit.DAGCircuit, BoolCircuit.DAGGate, BoolCircuit.GadgetCorrect, BoolCircuit.GateKind, BoolCircuit.rewriteGates, BoolCircuit.rewriteGates_inv}
Let $C$ be a circuit on $n$ inputs whose gates are in topological order and whose output
is a vertex of the circuit. Given natural number $c$, predicates $Q$ on gate kinds and
arities and $P$ on gates, and an emitter that is gadget-correct for these parameters,
assume every gate of $C$ satisfies $Q$ for its kind and arity. The rewritten circuit has
the gate list produced by processing the gates of $C$ from left to right, and its output
is the rewritten index corresponding to the original output, using $0$ if that index is
out of range.
\end{definition}

\begin{theorem}[Evaluation preserved by rewriting]
\lean{BoolCircuit.DAGCircuit.rewrite_eval}
\label{BoolCircuit.DAGCircuit.rewrite_eval}
\leanok
\uses{BoolCircuit.DAGCircuit.eval, BoolCircuit.DAGCircuit.rewrite, BoolCircuit.rewriteGates_inv}
\difficulty{1}
Let $C$ be a finite acyclic Boolean circuit with $n$ input vertices, let $Q$ be a
predicate on gate kinds and input counts, and let $P$ be a property of gates. Suppose a
rewriting procedure replaces each gate satisfying $Q$ by a gadget that extends the
existing acyclic gate list, preserves that gate's value for every input assignment, and
returns a valid output vertex. Require also that the gadget's gate list has length at
most the original list length plus the gate's number of inputs plus $2$, that its output
depth is at most $c$ plus the maximum depth of the gate's input vertices, and that every
added gate satisfies $P$. If every gate of $C$ satisfies $Q$, then for every input
assignment the rewritten circuit and $C$ have the same output value.
\end{theorem}

\begin{theorem}[Depth bound for circuit rewriting]
\lean{BoolCircuit.DAGCircuit.rewrite_depth_le}
\label{BoolCircuit.DAGCircuit.rewrite_depth_le}
\leanok
\uses{BoolCircuit.DAGCircuit.depth, BoolCircuit.DAGCircuit.rewrite, BoolCircuit.rewriteGates_inv}
\difficulty{1}
Let $C$ be a finite acyclic Boolean circuit with $n$ input vertices, let $Q$ be a
predicate on gate kinds and input counts, and let $P$ be a property of gates. Suppose a
rewriting procedure replaces each gate satisfying $Q$ by a gadget that extends the
existing acyclic gate list, preserves that gate's value for every input assignment, and
returns a valid output vertex. Require also that the gadget's gate list has length at
most the original list length plus the gate's number of inputs plus $2$, that its output
depth is at most $c$ plus the maximum depth of the gate's input vertices, and that every
added gate satisfies $P$. If every gate of $C$ satisfies $Q$, then the rewritten circuit
has depth at most $c$ times the depth of $C$.
\end{theorem}

\begin{theorem}[Property of rewritten gates]
\lean{BoolCircuit.DAGCircuit.rewrite_new_gates}
\label{BoolCircuit.DAGCircuit.rewrite_new_gates}
\leanok
\uses{BoolCircuit.DAGCircuit.rewrite, BoolCircuit.rewriteGates_inv}
\difficulty{1}
Let $C$ be a finite acyclic Boolean circuit with $n$ input vertices, let $Q$ be a
predicate on gate kinds and input counts, and let $P$ be a property of gates. Suppose a
rewriting procedure replaces each gate satisfying $Q$ by a gadget that extends the
existing acyclic gate list, preserves that gate's value for every input assignment, and
returns a valid output vertex. Require also that the gadget's gate list has length at
most the original list length plus the gate's number of inputs plus $2$, that its output
depth is at most $c$ plus the maximum depth of the gate's input vertices, and that every
added gate satisfies $P$. If every gate of $C$ satisfies $Q$, then every gate of the
rewritten circuit satisfies $P$.
\end{theorem}

\begin{theorem}[Gate count bound for circuit rewriting]
\lean{BoolCircuit.DAGCircuit.rewrite_length_le}
\label{BoolCircuit.DAGCircuit.rewrite_length_le}
\leanok
\uses{BoolCircuit.DAGCircuit.rewrite, BoolCircuit.rewriteGates_inv}
\difficulty{1}
Let $C$ be a finite acyclic Boolean circuit with $n$ input vertices, let $Q$ be a
predicate on gate kinds and input counts, and let $P$ be a property of gates. Suppose a
rewriting procedure replaces each gate satisfying $Q$ by a gadget that extends the
existing acyclic gate list, preserves that gate's value for every input assignment, and
returns a valid output vertex. Require also that the gadget's gate list has length at
most the original list length plus the gate's number of inputs plus $2$, that its output
depth is at most $c$ plus the maximum depth of the gate's input vertices, and that every
added gate satisfies $P$. If every gate of $C$ satisfies $Q$, then the rewritten circuit
has at most $\sum_{g\in G} (|\operatorname{args}(g)|+2)$ gates, where $G$ is the gate
list of $C$.
\end{theorem}

\begin{theorem}[Bound on total gate cost]
\lean{BoolCircuit.DAGCircuit.sum_cost_le}
\label{BoolCircuit.DAGCircuit.sum_cost_le}
\leanok
\uses{BoolCircuit.DAGCircuit, BoolCircuit.DAGCircuit.size}
\difficulty{4}
Let $C$ be a directed acyclic Boolean circuit with $n$ inputs, and let its size be $n$
plus its number of gates. If every gate of $C$ has at most $\mathrm{size}(C)$ arguments,
then the sum over all gates of their number of arguments plus two is at most the number
of gates times $\mathrm{size}(C)+2$.
\end{theorem}
